% answers_short.tex —— 2005 年·普通高等学校招生考试·北京卷·理科 · 简短答案速查
% 由 tools/make_short_answers.py 自 answers.tex 的 \gksolhead/\gkanswer 提取，
% 并与 2005/2005_数学真题_北京理科_答案要点.md 逐题比对后生成，勿手改（改源 MD 后重跑）。

\begin{gknotice}
\textbf{使用说明}：本节只给最终答案便于快速核对，完整推导见解析版。答案为编写者独立推导，\textbf{非官方标准答案}。
\end{gknotice}

\gkpart{选择题}
\begin{center}\begin{tabular}{ccccccccc}
\toprule
题号 & 1 & 2 & 3 & 4 & 5 & 6 & 7 & 8 \\
\midrule
答案 & C & B & C & B & D & C & A & A \\
\bottomrule
\end{tabular}\end{center}

\gkpart{填空题}
\begin{itemize}[leftmargin=2.6em]
  \item[\textbf{9.}] $a = \dfrac{8}{3}$。
  \item[\textbf{10.}] $-\dfrac{4}{3}$；$-\dfrac{1}{7}$。
  \item[\textbf{11.}] $-20$。
  \item[\textbf{12.}] $(1, \mathrm{e})$；$\mathrm{e}$。
  \item[\textbf{13.}] ②③。
  \item[\textbf{14.}] $65$；$20$。
\end{itemize}

\gkpart{解答题}
\begin{itemize}[leftmargin=2.6em]
  \item[\textbf{15.}] (1) 单调减区间为 $(-\infty, -1]$ 和 $[3, +\infty)$；(2) $-7$
  \item[\textbf{16.}] (1) 见解析版；(2) $90^\circ$（或 $\dfrac{\pi}{2}$）；(3) $\arccos\dfrac{\sqrt{15}}{5}$
  \item[\textbf{17.}] (1) $\xi$ 的概率分布列为 $P(\xi=0)=\dfrac{1}{8}, P(\xi=1)=\dfrac{3}{8}, P(\xi=2)=\dfrac{3}{8}, P(\xi=3)=\dfrac{1}{8}$；数学期望 $E\xi = \dfrac{3}{2}$；(2) $\dfrac{19}{27}$；(3) $\dfrac{1}{24}$
  \item[\textbf{18.}] (1) $W_1$ 表示为 $\begin{cases} x < 0, \\ kx < y < -kx \end{cases}$；$W_2$ 表示为 $\begin{cases} x > 0, \\ -kx < y < kx \end{cases}$；(2) $k^2 x^2 - y^2 = (k^2 + 1)d^2$；(3) 见解析版
  \item[\textbf{19.}] (1) $a_2 = a + \dfrac{1}{4}$，$a_3 = \dfrac{1}{2}a + \dfrac{1}{8}$；(2) 见解析版；(3) $2a - \dfrac{1}{2}$
  \item[\textbf{20.}] (1) 见解析版；(2) 见解析版；(3) $x_1 = 0.34, x_2 = 0.66, x_3 = 0.32$（或 $x_1 = 0.32, x_2 = 0.66, x_3 = 0.34$）
\end{itemize}
